\chapter{Выводы}
\label{ch:conclusions}

% В результате экспериментально были определены радиус линзы, создающей интерференционную картину в виде колец Ньютона, и разница длин волн зеленого и желтого компонентов света ртутной лампы. Радиус линзы оказался равным $(12,98 \pm 0,05) \ \text{мм}$. Полученное значение разности длин волн $(34,1 \pm 2,1) \ \text{нм}$ совпало с табличным $33 \ \text{нм}$ в пределах погрешности. Это означает применимость данного метода для точного определения разницы длин волн двух спектральных компонент света.

В результате было показано, что изучение интерференционной картины в виде колец Ньютона может быль использовано для определения радиуса кривизны линз и разницы длин волн двух узкополосных спектральных компонент света. Полученные результаты измерения на примере линейчатого спектра ртутной лампы для разности длин волн $(34,1 \pm 2,1) \ \text{нм}$ совпали с табличным значением $33,0 \ \text{нм}$ в пределах погрешности. Для исследуемой линзы радиус оказался равным $12,98 \ \text{мм}$ с погрешностью $0,05 \ \text{мм}$, что показало применимость исследования колец Ньютона для точного определения радиусов линз.